\begin{sol}[Rueda empujada por una fuerza horizontal en su centro]

\subsubsection*{Planteamiento general}

Tomamos $x$ horizontal (positivo en el sentido de $\vec{F}$) y $z$ saliendo de la página (convención antihoraria positiva). El momento de inercia de un aro delgado respecto a su centro es $I_G = mR^2$. Como $F$ y $m\vec{g}$ actúan sobre líneas que pasan por $G$ (o por su proyección vertical), solo la fricción $f$ (aplicada en $C$, a distancia $R$ bajo $G$) genera torque respecto a $G$:
$$\tau_z = R f$$

La condición de rodadura sin deslizar hacia la derecha requiere rotación horaria, es decir $\alpha_z = -a_G/R$.

\subsubsection*{Pregunta 1: Aceleración y fricción (rodadura sin deslizar)}

\textbf{Traslación:} $F + f = ma_G$

\textbf{Rotación:} $I_G\alpha_z = Rf \;\Rightarrow\; \alpha_z = \dfrac{Rf}{mR^2} = \dfrac{f}{mR}$

Usando $\alpha_z = -a_G/R$:
$$-\frac{a_G}{R} = \frac{f}{mR} \quad\Rightarrow\quad f = -ma_G$$

Sustituyendo en la ecuación de traslación:
$$F - ma_G = ma_G \quad\Rightarrow\quad \boxed{a_G = \frac{F}{2m}}$$

$$\boxed{f = -\frac{F}{2}}$$

El signo negativo indica que la fricción apunta en sentido contrario a $\vec{F}$, con magnitud $F/2$. Esto es consistente: al empujar exactamente por el centro, no se aplica torque directo, así que toda la aceleración angular necesaria para rodar debe generarla la fricción; al mismo tiempo, esa fricción reduce la fuerza neta de traslación, por lo cual $a_G = F/(2m)$ es la mitad de lo que sería sin rotación ($F/m$).

\subsubsection*{Pregunta 2: Fuerza máxima para rodadura sin deslizar}

En la dirección vertical, $N = mg$ (no hay aceleración vertical). La condición de no deslizamiento es $|f| \leq \mu_s N$:

$$\frac{F}{2} \leq \mu_s mg$$

$$\boxed{F_{\text{max}} = 2\mu_s mg}$$

\subsubsection*{Pregunta 3: Caso con deslizamiento ($F > F_{\text{max}}$)}

Una vez que el aro desliza, la fricción pasa a ser cinética, con magnitud fija $f_k = \mu_k mg$ (usamos $\mu_k$ para el coeficiente cinético; si no se especifica, se puede aproximar $\mu_k \approx \mu_s$), dirigida siempre oponiéndose al deslizamiento relativo en $C$ (hacia atrás, igual que en el caso anterior).

\textbf{Traslación:}
$$F - \mu_k mg = ma_G \quad\Rightarrow\quad \boxed{a_G = \frac{F}{m} - \mu_k g}$$

\textbf{Rotación:}
$$I_G\alpha_z = -R\mu_k mg \quad\Rightarrow\quad \alpha_z = -\frac{\mu_k mgR}{mR^2} = -\frac{\mu_k g}{R}$$

$$\boxed{|\alpha| = \frac{\mu_k g}{R}}$$

\paragraph{Interpretación física} A diferencia del caso de rodadura pura (donde $f$ se ajusta exactamente al valor necesario), al deslizar la fricción queda "saturada" en su valor máximo $\mu_k mg$ y ya no puede seguir el ritmo de $F$: por eso $a_G$ crece indefinidamente con $F$, mientras que $\alpha$ permanece constante (fijado únicamente por $\mu_k$, $g$ y $R$), y el aro empieza a patinar cada vez más respecto al suelo.

\end{sol}
